\section{Discussion and Limitations}
\label{sec:discussion}

\noindent\textbf{Claim language.}
We deliberately use the weaker framing ``observed effective levers
under single-lever interventions'' rather than ``independent causal
levers.'' The independence assumption—that $\delta_i$ and $\delta_j$
are additive—is not tested by our design. Levers in $E(x, a)$ are
tested separately against the unedited original; their joint effect
and potential interactions are left for future work requiring
factorial intervention designs.

\noindent\textbf{Generator leakage.}
Because our implementation relies on prompt-only editing rather than
mask-constrained generation, generators can introduce residual texture
and style changes outside the intended support.
The 2AFC measures the total image difference, not the isolated
effect of the intended edit. This is why we position the work as
interventional attribution rather than structural causal inference.

\noindent\textbf{Ontology scope.}
We measure $E(x, a)$ only with respect to the closed candidate set
defined by our ontology and scene-grounding procedure. Some potential
street-view cues may be outside that scope or excluded because they are
not judged safely manipulable under the prompt-only protocol.
$|E|{=}0$ may reflect this scope limitation rather than a
perception ceiling, particularly for high-baseline scenes.

\noindent\textbf{Grounding bias.}
Scene support grounded by planners or annotators may favor large,
salient, or easy-to-describe elements. Small but perceptually active
cues (a barely visible security camera, a subtle entry light) may be
absent from $\mathcal{L}(x)$.
Our multi-lever counts are therefore likely underestimates.
